\chapter{Performance Evaluation and Discussion}
\label{chap:performance_evaluation}

This chapter presents the quantitative evaluation of the deployed vineyard disease detection system on the ESP32-S3 platform. Section~\ref{sec:experimental_setup} describes the experimental setup. Section~\ref{sec:latency_analysis} provides a detailed latency analysis of each inference pipeline stage. Section~\ref{sec:accuracy_validation} validates model accuracy after quantization and deployment. Section~\ref{sec:memory_utilization} examines memory footprint and runtime stability. Section~\ref{sec:power_consumption} characterizes power consumption and duty-cycle energy efficiency. Finally, Section~\ref{sec:discussion} discusses the trade-offs inherent in the system design and acknowledges current limitations.

All measurements reported in this chapter are obtained from the production firmware (v29) running on the Freenove ESP32-S3-WROOM development board~\cite{freenove2024esp32s3} with OV3660 camera module, using serial monitor output at 115,200 baud for timing and memory profiling data.

\section{Experimental Setup}
\label{sec:experimental_setup}

The experiment was conducted in a controlled laboratory environment to validate the complete inference pipeline on the target hardware. The production firmware (v29) was flashed onto the Freenove ESP32-S3-WROOM development board via USB, which also served as the power source (5~V USB, regulated to 3.3~V internally) and as the serial monitoring interface at 115,200~baud. The OV3660 camera module was connected to the board's camera ribbon connector.

A printed grapevine banner depicting vine leaves was used as the test subject. The camera was directed at the banner by hand to simulate field-like capture conditions. The detection confidence threshold was set to $\tau = 0.25$ and the maximum detection count was limited to 10 leaves per frame to ensure sufficient detections for testing the full pipeline. Each detected leaf region was cropped and passed through the MobileNetV2 classifier, with per-leaf disease predictions aggregated into a frame-level diagnosis. All inference results, timing breakdowns, and memory statistics were printed to the serial monitor for analysis.

\begin{itemize}
    \item \textbf{Hardware:} Freenove ESP32-S3-WROOM development board (dual-core Xtensa LX7 at 240~MHz, 8~MB PSRAM, 16~MB flash) with OV3660 camera module.
    \item \textbf{Firmware:} Production firmware v29 implementing the complete pipeline: camera capture, JPEG decode, ESPDet-Pico leaf detection, MobileNetV2 disease classification, and hybrid aggregation.
    \item \textbf{Test input:} Printed grapevine banner with handheld camera positioning, providing a controlled and repeatable input for end-to-end pipeline validation.
    \item \textbf{Timing instrumentation:} Microsecond-resolution timers embedded in the firmware at each pipeline stage, output to the serial monitor.
    \item \textbf{Memory profiling:} Free PSRAM readings logged at each initialization stage and before/after each inference cycle to verify memory stability.
    \item \textbf{Power characterization:} Active-mode current derived from the ESP32-S3 datasheet specification for 240~MHz dual-core operation with octal PSRAM active~\cite{esp32s3_datasheet}; deep sleep current (8~$\mu$A) taken from the datasheet specification for RTC-timer wakeup mode.
\end{itemize}


\section{Inference Latency Analysis}
\label{sec:latency_analysis}

\subsection*{Per-Stage Latency Breakdown}
\label{subsec:per_stage_latency}

The inference pipeline is instrumented with microsecond-resolution timing measurements at each processing stage. Table~\ref{tab:per_stage_latency} presents the measured latency for each stage during a representative inference cycle capturing a grapevine image with 10 detected leaves.

\begin{table}[htbp]
\centering
\caption{Per-stage latency breakdown for the inference pipeline (10 detected leaves)}
\label{tab:per_stage_latency}
\begin{tabular}{lrc}
\toprule
\textbf{Pipeline Stage} & \textbf{Latency (ms)} & \textbf{Percentage} \\
\midrule
Camera capture (JPEG)                                    & 78     & 1.2\%   \\
JPEG decode (RGB888)                                     & 51     & 0.8\%   \\
Leaf detection (incl.\ letterbox preprocessing)           & 361    & 5.7\%   \\
Disease classification ($\times$10 leaves, incl.\ crop)   & 5,886  & 92.3\%  \\
\midrule
\textbf{Total end-to-end}                                  & \textbf{6,376}  & \textbf{100\%} \\
\bottomrule
\end{tabular}
\end{table}

MobileNetV2 classification dominates the pipeline at 92.3\% of total latency. The classification time of 5,886~ms for 10 leaves corresponds to approximately 589~ms per leaf, which decomposes into crop and resize preprocessing ($\sim$1--2~ms), tensor setup ($\sim$0.1~ms), the INT8 forward pass ($\sim$550~ms), and softmax post-processing ($\sim$0.1~ms). The forward pass constitutes approximately 93\% of the per-leaf classification time and is dominated by depthwise separable convolution operations executing in INT8 arithmetic on the Xtensa LX7 SIMD units.

The detection stage accounts for 5.7\% of total latency despite processing the full $416 \times 320$ input. This time includes letterbox preprocessing ($\sim$22~ms), the ESPDet-Pico INT8 forward pass ($\sim$340~ms), and non-maximum suppression post-processing. The favorable detection speed relative to classification reflects ESPDet-Pico's substantially smaller parameter count (360K vs $\sim$2.2M for MobileNetV2) and the computational efficiency of the anchor-free single-pass architecture.

Image acquisition stages (capture and JPEG decode) together contribute only 2.0\% of total latency, confirming that neural network inference --- not I/O or preprocessing --- constitutes the computational bottleneck in embedded edge AI deployments.

\subsection*{End-to-End Pipeline Latency and Throughput}
\label{subsec:e2e_latency}

The end-to-end latency scales linearly with the number of detected leaves, as classification is applied sequentially to each detected region:

\begin{equation}
T_{\text{total}} = T_{\text{capture}} + T_{\text{decode}} + T_{\text{detect}} + N \cdot (T_{\text{crop}} + T_{\text{classify}}) + T_{\text{aggregate}}
\end{equation}

where $N$ denotes the number of detected leaves and $(T_{\text{crop}} + T_{\text{classify}}) \approx 589$~ms per leaf (of which the INT8 forward pass accounts for $\sim$550~ms). This linear relationship yields the following throughput estimates:

\begin{table}[htbp]
\centering
\caption{End-to-end latency and throughput versus number of detected leaves}
\label{tab:latency_vs_leaves}
\begin{tabular}{ccc}
\toprule
\textbf{Leaves Detected} & \textbf{Estimated Latency (ms)} & \textbf{Throughput (FPS)} \\
\midrule
1   & $\sim$1,080   & 0.93  \\
3   & $\sim$2,260   & 0.44  \\
5   & $\sim$3,440   & 0.29  \\
7   & $\sim$4,610   & 0.22  \\
10  & 6,376         & 0.16  \\
\bottomrule
\end{tabular}
\end{table}

At the maximum detection count of 10 leaves per frame, the system achieves 0.16 FPS, corresponding to one complete inference cycle every 6.4 seconds. For the target application of periodic vineyard monitoring --- where disease symptoms evolve over days rather than hours --- the system performs four inference cycles per day at 6-hour intervals. The resulting duty cycle is negligible (6.4~s active out of 21,600~s per cycle $\approx$ 0.03\%), ensuring that latency imposes no operational constraint on the monitoring schedule.


\section{Model Accuracy Validation}
\label{sec:accuracy_validation}

This section evaluates the accuracy of both models in their quantized INT8 form deployed on the ESP32-S3, comparing against the FP32 training baselines established in Chapter~\ref{chap:system_architecture}.

\subsection*{ESPDet-Pico Detection Accuracy}
\label{subsec:espdet_accuracy}

Table~\ref{tab:espdet_accuracy_comparison} compares the detection performance of the FP32 training model against the post-quantization INT8 model evaluated on the held-out test set (103 images).

\begin{table}[htbp]
\centering
\caption{ESPDet-Pico detection accuracy: FP32 vs INT8 quantized}
\label{tab:espdet_accuracy_comparison}
\begin{tabular}{lccc}
\toprule
\textbf{Metric} & \textbf{FP32} & \textbf{INT8} & \textbf{Degradation} \\
\midrule
mAP@50 (\%)         & 56.42 & 55.8  & $-$0.62 pp \\
mAP@50-95 (\%)      & 32.29 & 31.5  & $-$0.79 pp \\
Precision (\%)      & 62.77 & 61.9  & $-$0.87 pp \\
Recall (\%)         & 54.17 & 53.5  & $-$0.67 pp \\
\midrule
Model size          & 1.06~MB & 478~KB & $-$55\% \\
\bottomrule
\end{tabular}
\end{table}

INT8 accuracy values are obtained by evaluating the quantized model on the same held-out test set used for FP32 validation, using the ESP-PPQ framework's built-in accuracy evaluation pipeline. The quantization introduces less than 1 percentage point degradation across all detection metrics. The mAP@50 decrease from 56.42\% to 55.8\% represents a 1.1\% relative accuracy loss, well within acceptable bounds for the leaf localization task where the detector serves as a region proposal mechanism for the downstream classifier. The minimal degradation reflects the effectiveness of the layerwise equalization strategy (Section~\ref{subsec:quantization_accuracy}) in preserving detection performance despite aggressive 32-bit to 8-bit precision reduction.

The quantized model retains sufficient detection capability for the downstream classification pipeline: at the production confidence threshold ($\tau = 0.25$), the detector reliably localizes leaves in vineyard imagery, providing region proposals for disease classification.

\subsection*{MobileNetV2 Classification Accuracy}
\label{subsec:mobilenet_accuracy}

Table~\ref{tab:mobilenet_accuracy_comparison} presents the classification accuracy comparison between the FP32 and INT8 models on the held-out test set (814 images).

\begin{table}[htbp]
\centering
\caption{MobileNetV2 classification accuracy: FP32 vs INT8 quantized}
\label{tab:mobilenet_accuracy_comparison}
\begin{tabular}{lccc}
\toprule
\textbf{Metric} & \textbf{FP32} & \textbf{INT8} & \textbf{Degradation} \\
\midrule
Test Accuracy (\%)    & 99.80 & $>$99.5 & $<$0.3 pp \\
Macro F1-score        & 0.998 & $>$0.995 & $<$0.003 \\
\midrule
Model size (ONNX)     & 8.48~MB & 2.27~MB & $-$73.2\% \\
\bottomrule
\end{tabular}
\end{table}

Evaluation of the INT8 quantized model on the same held-out test set used for FP32 validation indicates accuracy above 99.5\%, with less than 0.3 percentage point degradation from the FP32 baseline. The robustness to quantization is attributed to three factors: (1) batch normalization layers in MobileNetV2's inverted residual blocks stabilize activation distributions, reducing sensitivity to precision reduction; (2) the layerwise equalization technique redistributes weight magnitudes to minimize per-layer quantization error; and (3) the moderate complexity of the four-class classification task (compared to ImageNet's 1000 classes) provides inherent margin for precision loss.

\subsection*{Aggregation Strategy Evaluation}
\label{subsec:aggregation_evaluation}

The three multi-leaf aggregation strategies (Section~\ref{subsec:aggregation_strategies}) were evaluated on the system integration test using a printed grapevine banner with OV3660 camera capture. Table~\ref{tab:aggregation_sample} presents a representative frame-level aggregation result from the hybrid spatial-quality weighting strategy with 10 detected leaves.

\begin{table}[htbp]
\centering
\caption{Sample aggregation output: hybrid spatial-quality weighting (10 leaves)}
\label{tab:aggregation_sample}
\begin{tabular}{lc}
\toprule
\textbf{Aggregated Class Score} & \textbf{Probability} \\
\midrule
Black Rot       & 9.93\% \\
Esca            & 7.54\%  \\
\textbf{Healthy}& \textbf{62.58\%} \\
Leaf Blight     & 19.94\% \\
\midrule
\textbf{Frame diagnosis:} & Healthy (62.6\% confidence) \\
\bottomrule
\end{tabular}
\end{table}

The hybrid aggregation strategy assigns per-leaf weights based on the product of detection confidence, classification certainty (1 $-$ normalized entropy), and bounding box quality. Observed leaf weights range from 0.0123 (low-quality edge detection with high entropy, 12.3\% certainty) to 0.0695 (well-detected leaf with moderate certainty, 33.9\%). Leaves with high classification entropy (indicating ambiguous predictions near class boundaries) are effectively downweighted, reducing the influence of unreliable classifications on the frame-level diagnosis.

The aggregation computation adds less than 2~ms to the pipeline latency, introducing negligible overhead. A comprehensive comparative evaluation of the three aggregation strategies on field test data is planned as part of the on-field validation described in Section~\ref{subsec:limitations}.


\section{Memory Utilization and Runtime Stability}
\label{sec:memory_utilization}

\subsection*{Memory Footprint Analysis}
\label{subsec:memory_footprint}

Table~\ref{tab:memory_footprint} presents the PSRAM allocation breakdown measured from the boot log serial output of the production firmware, showing free PSRAM at each initialization stage.

\begin{table}[htbp]
\centering
\caption{PSRAM allocation breakdown during system initialization}
\label{tab:memory_footprint}
\begin{tabular}{lrrr}
\toprule
\textbf{Initialization Stage} & \textbf{Free PSRAM (KB)} & \textbf{Consumed (KB)} & \textbf{Time (ms)} \\
\midrule
System boot (heap init)            & 8,190  & ---      & ---  \\
After ESPDet-Pico init             & 7,095  & 1,095    & 144  \\
After MobileNetV2 init             & 4,207  & 2,888    & 348  \\
After camera init + runtime        & ---    & $\sim$297 & ---  \\
During inference (steady state)    & 3,910  & ---      & ---  \\
\bottomrule
\end{tabular}
\end{table}

The total PSRAM consumption is approximately 4,280~KB (52.2\% of 8~MB), leaving 3,910~KB free during steady-state inference. This memory budget includes:

\begin{itemize}
    \item \textbf{ESPDet-Pico activation buffers:} $\sim$1,095~KB for detection model activation tensors and internal buffers.
    \item \textbf{MobileNetV2 activation buffers:} $\sim$2,839~KB for classification model activation tensors plus the 49~KB persistent input buffer.
    \item \textbf{Camera DMA buffers:} $2 \times 60$~KB = 120~KB for double-buffered JPEG frame capture in PSRAM.
    \item \textbf{Driver and runtime overhead:} $\sim$177~KB for camera driver allocations, heap metadata, and memory alignment.
    \item \textbf{JPEG decode buffer:} $\sim$900~KB (dynamically allocated and freed per frame).
\end{itemize}

The 3,910~KB free PSRAM during inference provides sufficient margin for the temporary JPEG decode buffer allocation ($\sim$900~KB) and other runtime overhead, with approximately 3~MB remaining as safety margin against unexpected allocations or memory alignment overhead.

Flash memory allocation is dominated by the firmware binary (4,356,032 bytes = 4.15~MB of 5~MB factory partition), of which 2,971,836 bytes (2.83~MB) is memory-mapped model data and 1,268,968 bytes (1.21~MB) is application code. Internal SRAM (274~KB of 512~KB available at boot for dynamic allocation) is used by the ESP-DL runtime for weight tile caching (up to 64~KB) and FreeRTOS task stacks.

\subsection*{Long-Term Stability Validation}
\label{subsec:long_term_stability}

The single-allocation zero-copy buffer strategy (Section~\ref{subsec:buffer_allocation}) is validated through continuous operation testing. The system monitors free PSRAM before and after each inference cycle, verifying that the post-inference free memory returns to the same baseline value (3,910~KB $\pm$ 1~KB) after every cycle. This confirms the absence of memory leakage from unfreed dynamic allocations or fragmentation-induced memory loss.

The deterministic memory behavior stems from the design decision to eliminate all dynamic allocations within the inference loop: the JPEG decode buffer is the only runtime allocation, and its allocation-free pattern is verified through explicit size assertions. All other buffers (classification input, model activations, camera framebuffers) are allocated once during initialization and reused across all subsequent inference cycles without reallocation.


\section{Power Consumption analysis}
\label{sec:power_consumption}

\subsection{Active and Sleep Power Measurements}
\label{subsec:power_measurements}

Power consumption is characterized for two operational states. Active-mode power is derived from the ESP32-S3 datasheet maximum current specification~\cite{esp32s3_datasheet} for 240~MHz dual-core operation with peripherals active; deep sleep current is taken from the datasheet specification for RTC-timer wakeup mode:

\begin{itemize}
    \item \textbf{Active mode (inference):} During neural network inference at 240~MHz with PSRAM access, camera operation, and flash reads, the system draws approximately 74~mA at 3.3~V (244~mW). This represents the ESP32-S3's maximum operating current under full dual-core computational load with octal PSRAM active at 80~MHz. Power consumption during inference is dominated by the Xtensa LX7 cores executing INT8 SIMD operations and the octal PSRAM interface transferring activation data at each layer boundary.
    \item \textbf{Deep sleep mode:} Between inference cycles, the system enters deep sleep with all peripherals powered down (camera, PSRAM controller, flash interface). Only the RTC domain remains active for timed wakeup. The deep sleep current is 8~$\mu$A (26~$\mu$W at 3.3~V), per the ESP32-S3 datasheet specification for deep sleep with RTC timer wakeup enabled.
\end{itemize}

\subsection{Duty-Cycle Energy Analysis}
\label{subsec:duty_cycle}

Since grapevine foliar diseases such as Black Rot (\textit{Guignardia bidwellii}) and Esca develop visible symptoms over days to weeks, the system's condition does not change appreciably between consecutive captures at sub-hour intervals. Accordingly, the deployed monitoring schedule performs four inference cycles per day at 6-hour intervals (e.g., 06:00, 12:00, 18:00, 00:00), capturing disease state across the diurnal lighting cycle while spending the vast majority of time in deep sleep. Each inference cycle completes in approximately 6.4 seconds for a 10-leaf frame. The per-cycle duty cycle is:

\begin{equation}
D = \frac{T_{\text{active}}}{T_{\text{cycle}}} = \frac{6.4~\text{s}}{21{,}600~\text{s}} \approx 0.030\%
\end{equation}

The average power consumption over one complete duty cycle combines the active and sleep contributions:

\begin{equation}
P_{\text{avg}} = D \cdot P_{\text{active}} + (1 - D) \cdot P_{\text{sleep}} = 2.96 \times 10^{-4} \times 244~\text{mW} + 0.9997 \times 0.026~\text{mW} \approx 0.098~\text{mW}
\label{eq:avg_power}
\end{equation}

At an average consumption of 98~$\mu$W, the system's daily energy requirement (four cycles) is approximately 2.4~mWh. A standard 3.7~V, 2000~mAh lithium-polymer battery (7,400~mWh capacity) could theoretically sustain the system for over 8 years without recharging, far exceeding the practical battery shelf life. In practice, battery self-discharge ($\sim$3--5\% per month for LiPo cells) becomes the dominant energy loss mechanism, yielding realistic autonomous operation of 6--12 months depending on environmental temperature. With a modest solar panel (e.g., 1~W rated capacity receiving 4 effective sun-hours per day), the system achieves indefinite energy-autonomous operation throughout the growing season and beyond.

The 6-hour monitoring interval is well-matched to the temporal dynamics of grapevine disease progression. The capture interval is configurable via the \texttt{CAPTURE\_INTERVAL\_SEC} firmware parameter, allowing operators to increase monitoring frequency during periods of high disease pressure (e.g., humid conditions favoring fungal growth) at the cost of proportionally higher energy consumption.

\section{Discussion}
\label{sec:discussion}

\subsection{Trade-offs Between Accuracy, Latency, and Power}
\label{subsec:tradeoffs}

The deployed system embodies a three-way engineering trade-off between accuracy, inference latency, and power consumption:

\textbf{Accuracy vs. Model Size:} ESPDet-Pico sacrifices 8.6\% relative mAP@50 compared to YOLO11n (56.42\% vs 61.75\%) in exchange for an 86\% parameter reduction that enables in-memory deployment. MobileNetV2 at $128 \times 128$ input resolution trades potential accuracy gains from higher resolution ($224 \times 224$) for a 3$\times$ reduction in activation memory, enabling co-deployment with the detection model within the 8~MB PSRAM budget. Despite these compromises, the classification model achieves 99.80\% accuracy on the controlled test set, indicating that the accuracy trade-offs are well justified for this task complexity.

\textbf{Latency vs. Accuracy:} The choice of nearest-neighbor interpolation over bilinear for crop resizing sacrifices image quality for 5--8$\times$ speedup per crop operation. INT8 quantization provides $\sim$4$\times$ computational speedup via SIMD utilization at the cost of marginal accuracy degradation ($<$0.3 pp for classification). The maximum detection limit of 10 leaves per frame bounds worst-case latency at 6.4 seconds while potentially missing additional leaves in dense canopy regions.

\textbf{Power vs. Throughput:} The 6-hour capture interval (four times daily) trades temporal resolution for extreme energy efficiency, reducing average power consumption to 98~$\mu$W --- a level at which battery self-discharge dominates over system consumption. For applications requiring higher monitoring frequency (e.g., hourly captures during humid conditions favoring rapid fungal onset), the duty cycle can be increased at the cost of proportionally higher energy consumption, though even at 5-minute intervals the system remains well within solar-harvesting feasibility.

\subsection{Limitations and Deployment Considerations}
\label{subsec:limitations}

The current evaluation has several limitations that should be acknowledged:

\textbf{Laboratory vs. field accuracy:} All accuracy metrics reported in Sections~\ref{subsec:espdet_accuracy} and~\ref{subsec:mobilenet_accuracy} are obtained from controlled test sets (PlantVillage dataset for classification, Roboflow-annotated images for detection). The system integration test described in Section~\ref{subsec:aggregation_evaluation} uses a printed grapevine banner to validate the inference pipeline end-to-end, but does not constitute field-level accuracy validation. On-field testing with live grapevine canopy under natural environmental conditions (variable lighting, wind-induced leaf motion, rain droplets, dust accumulation on camera lens) remains a planned future activity, pending vineyard preparation for the upcoming growing season.

\textbf{Environmental robustness:} The OV3660 camera's automatic exposure and white balance provide basic adaptation to lighting variations. However, extreme conditions --- direct sunlight causing sensor saturation, heavy shadow creating underexposure, or fog/rain degrading image clarity --- may adversely affect both detection recall and classification accuracy. The impact of these environmental factors on system performance has not yet been quantified.

\textbf{Disease progression monitoring:} The current system provides single-frame disease classification without temporal tracking across consecutive capture cycles. Implementing change detection or disease progression monitoring would require storing historical predictions in NVS (non-volatile storage) and comparing across frames, which is feasible within the current firmware architecture but not yet implemented.

\textbf{Communication integration:} The deployed prototype outputs diagnostic results via serial monitor. Integration with LoRaWAN~\cite{adelantado2017understanding} for wireless transmission of disease alerts to a central gateway is architecturally planned (the system architecture in Chapter~\ref{chap:system_architecture} shows LoRaWAN integration) but not yet implemented in the current firmware. The 4.15~MB firmware binary leaves approximately 850~KB of the 5~MB factory partition for LoRaWAN stack integration.

\textbf{Aggregation strategy validation:} The comparative evaluation of the three aggregation strategies (baseline, entropy-based, hybrid) requires on-field data with known ground-truth disease labels across multiple leaves per frame. This validation is contingent on the vineyard field trials described above and is planned as a priority evaluation task for the next growing season.

Despite these limitations, the results presented in this chapter demonstrate that the complete inference pipeline --- from camera capture through multi-leaf aggregation --- executes reliably on a microcontroller-class device, achieving classification accuracy consistent with the FP32 training baselines while operating within the strict memory and power constraints of the ESP32-S3 platform. The implications of these findings and directions for future work are discussed in Chapter~\ref{chap:conclusion}.
